% =============================================================================
% CHƯƠNG 1 — XÁC ĐỊNH VẤN ĐỀ NGHIÊN CỨU
% Bản thảo tiếng Việt để duyệt nội dung trước khi ghép vào luận văn UEH.
% Nguồn nội dung chính: Main GBM_HESTON_SBTS_final (1).tex và các notebook 25_*.
% =============================================================================

\chapter{XÁC ĐỊNH VẤN ĐỀ NGHIÊN CỨU}
\label{chap:xac-dinh-van-de}

\section{Bối cảnh nghiên cứu}
\label{sec:boi-canh-nghien-cuu}

Quyền chọn đa tài sản (\textit{rainbow options}) là nhóm sản phẩm phái sinh có
giá trị thanh toán phụ thuộc đồng thời vào diễn biến của nhiều tài sản cơ sở.
Khác với quyền chọn đơn tài sản, cấu trúc thanh toán của quyền chọn rainbow có
thể sử dụng trung bình của một rổ tài sản, giá trị lớn nhất hoặc nhỏ nhất trong
rổ, hoặc kết hợp các đặc điểm này với phép lấy trung bình theo thời gian. Vì
vậy, rủi ro của hợp đồng không chỉ phụ thuộc vào phân phối biên của từng tài
sản mà còn phụ thuộc vào cấu trúc phụ thuộc chéo, sự đồng biến trong các giai
đoạn căng thẳng và toàn bộ quỹ đạo giá trong thời hạn hợp đồng
\parencite{Johnson1987,Wystup2006}.

Nghiên cứu xem xét hai hợp đồng đại diện. Hợp đồng thứ nhất là quyền chọn mua
Asian trên rổ tài sản (\textit{basket Asian call}); khoản thanh toán phụ thuộc
vào mức giá bình quân theo thời gian của cả rổ so với giá thực hiện. Hợp đồng
thứ hai là quyền chọn bán Asian trên tài sản kém nhất
(\textit{Asian worst-of put}); khoản thanh toán phụ thuộc vào tài sản có mức
giá bình quân thấp nhất trong rổ. Nói cách khác, \textit{basket Asian call}
phản ánh diễn biến chung của cả rổ, còn \textit{Asian worst-of put} tập trung
vào rủi ro đến từ thành phần suy giảm mạnh nhất. Nhờ đó, hai hợp đồng cho phép
đánh giá bộ sinh quỹ đạo ở cả điều kiện thị trường thuận lợi lẫn những tình
huống một hoặc nhiều tài sản diễn biến bất lợi.

Những khó khăn trong phòng ngừa rủi ro có thể được hình dung rõ hơn khi so
sánh với trường hợp Black--Scholes quen thuộc. Đối với quyền chọn châu Âu trên
một tài sản không trả cổ tức, mô hình Black--Scholes cho phép xác định giá
quyền chọn mua và quyền chọn bán bằng các công thức dạng đóng
\parencite{BlackScholes1973}:
\begin{align}
C_0 &= S_0\Phi(d_1)-Ke^{-rT}\Phi(d_2),\\
P_0 &= Ke^{-rT}\Phi(-d_2)-S_0\Phi(-d_1),
\end{align}
trong đó
\begin{equation}
d_1=\frac{\ln(S_0/K)+(r+\sigma^2/2)T}{\sigma\sqrt{T}},
\qquad
d_2=d_1-\sigma\sqrt{T}.
\end{equation}
Từ đó, các độ nhạy phục vụ phòng ngừa rủi ro cũng có thể được tính trực tiếp;
chẳng hạn, Delta của quyền chọn mua là $\Phi(d_1)$. Tuy nhiên, đối với các
hợp đồng được xem xét trong nghiên cứu này, giá trị thanh toán còn phụ thuộc
vào giá bình quân theo thời gian và, trong trường hợp quyền chọn
\textit{worst-of}, vào tài sản có kết quả thấp nhất trong rổ. Vì vậy, nhìn
chung không tồn tại những công thức đơn giản tương tự Black--Scholes. Ngay cả
quyền chọn Asian đơn tài sản cũng thường phải được định giá bằng phương pháp
xấp xỉ hoặc mô phỏng số \parencite{KemnaVorst1990}.

Khó khăn thứ hai xuất phát từ sự phụ thuộc giữa các tài sản. Hiệu quả phòng
ngừa rủi ro không chỉ phụ thuộc vào diễn biến riêng của từng tài sản mà còn
phụ thuộc vào cách chúng cùng biến động. Mô hình chuyển động Brown hình học đa
biến (\textsc{GBM}) giả định lợi suất log tuân theo phân phối Gaussian và mối
tương quan giữa các tài sản được biểu diễn bằng một ma trận tương quan cố
định. Mô hình Heston cho phép phương sai biến động ngẫu nhiên
\parencite{Heston1993}, nhưng việc mở rộng và hiệu chỉnh mô hình cho nhiều tài
sản vẫn cần thêm các giả định tham số để kiểm soát độ phức tạp
\parencite{Dimitroff2011}.

Khó khăn thứ ba đến từ chi phí giao dịch và việc tái cân bằng danh mục theo
các thời điểm rời rạc. Trong mô hình Black--Scholes lý tưởng, danh mục phòng
ngừa được giả định có thể điều chỉnh liên tục và không phát sinh chi phí.
Trong thực tế, mỗi lần giao dịch đều làm phát sinh chi phí, nên chiến lược sao
chép liên tục không thể được thực hiện chính xác như trong mô hình lý thuyết
\parencite{BlackScholes1973,Leland1985}.

Phòng ngừa rủi ro sâu (\textit{deep hedging}) cung cấp một hướng tiếp cận phù
hợp cho ba khó khăn trên. Thay vì suy ra chiến lược phòng ngừa từ một công
thức định giá đóng, phương pháp này tham số hóa chiến lược bằng mạng nơ-ron và
huấn luyện trực tiếp trên các quỹ đạo mô phỏng nhằm tối thiểu hóa một hàm rủi
ro của sai số phòng ngừa \parencite{Buehler2019}. Cách tiếp cận này cho phép
kết hợp quyền chọn phụ thuộc đường đi, nhiều tài sản cơ sở, tái cân bằng rời
rạc và chi phí giao dịch trong cùng một bài toán tối ưu.

Tuy nhiên, ưu điểm của deep hedging không loại bỏ rủi ro mô hình. Mạng phòng
ngừa học từ các quỹ đạo do một bộ sinh dữ liệu tạo ra; vì thế, phân phối của
dữ liệu huấn luyện trở thành một thành phần quyết định của chiến lược. Một
mạng nơ-ron được tối ưu tốt trên các quỹ đạo không đại diện cho thị trường vẫn
có thể tạo ra chiến lược kém hiệu quả khi áp dụng trên dữ liệu thực. Do đó,
việc lựa chọn bộ sinh quỹ đạo không chỉ là một bước kỹ thuật đầu vào mà là một
quyết định có ý nghĩa kinh tế và quản trị rủi ro.

Trong bối cảnh đó, mô hình chuỗi thời gian dựa trên cầu Schrödinger
(\textit{Schrödinger Bridge Time Series}, viết tắt là \textsc{SBTS}) là một
đối chứng dựa trên dữ liệu đáng chú ý. Thay vì áp đặt trước một dạng khuếch tán
tham số, \textsc{SBTS} xây dựng quỹ đạo từ phân phối thực nghiệm và cấu trúc
chuyển tiếp của dữ liệu lịch sử, trên nền tảng của bài toán cực tiểu hóa entropy
tương đối của cầu Schrödinger \parencite{Hamdouche2023}. Trong trường hợp đa
biến, thông tin trạng thái bao gồm đồng thời lợi suất của nhiều tài sản và độ
trễ Markov được lựa chọn từ dữ liệu; bandwidth và bậc nhớ được hiệu chỉnh qua
sai số dự báo ngoài mẫu \parencite{Alouadi2025}. Đặc điểm này tạo ra kỳ vọng
rằng \textsc{SBTS} có thể bảo toàn tốt hơn các quan hệ phi tuyến và phụ thuộc
đuôi trong dữ liệu tài chính đa biến.

Từ các lập luận trên, câu hỏi trung tâm không phải là bộ sinh nào tái tạo một
chỉ tiêu thống kê riêng lẻ tốt nhất, mà là: \emph{khác biệt giữa các bộ sinh có
chuyển hóa thành khác biệt có ý nghĩa trong một quyết định kinh tế hạ nguồn hay
không?} Nghiên cứu sử dụng bài toán phòng ngừa rủi ro sâu cho quyền chọn rainbow
dưới chi phí giao dịch để trả lời câu hỏi này.

\section{Vấn đề nghiên cứu và khoảng trống nghiên cứu}
\label{sec:van-de-khoang-trong}

\subsection{Vấn đề nghiên cứu}
\label{subsec:van-de-nghien-cuu}

Hiệu quả của một chiến lược deep hedging là kết quả đồng thời của ba thành
phần: phân phối quỹ đạo huấn luyện, kiến trúc và quy trình tối ưu của mạng
phòng ngừa, và chế độ thị trường dùng để đánh giá. Nếu thay đổi nhiều thành
phần cùng lúc, chênh lệch kết quả không thể được quy cho bộ sinh dữ liệu. Vì
vậy, nghiên cứu xây dựng một thiết kế so sánh có kiểm soát, trong đó
\textsc{GBM}, Heston và \textsc{SBTS} chỉ khác nhau ở phân phối quỹ đạo đầu
vào; các yếu tố còn lại gồm dữ liệu hiệu chỉnh, số quỹ đạo, kiến trúc mạng,
chi phí giao dịch, hợp đồng, strike, seed và quy trình huấn luyện được giữ
thống nhất.

Vấn đề nghiên cứu được xác định là đánh giá liệu bộ sinh đa biến dựa trên cầu
Schrödinger có tạo ra chiến lược phòng ngừa rủi ro tốt hơn các bộ sinh tham số
hay không, khi cả ba được đặt trong cùng một khuôn khổ deep hedging và được
đánh giá trên dữ liệu thị trường thực qua nhiều chế độ. Khái niệm ``tốt hơn''
được xem xét trên ba phương diện:

\begin{enumerate}
    \item \textbf{Hiệu quả phòng ngừa:} mức phân tán và rủi ro đuôi của sai số
    phòng ngừa sau khi đã tính chi phí giao dịch;
    \item \textbf{Hiệu quả vốn:} mức tiền mặt ban đầu $V_0^{*}$ được học trong
    giai đoạn tối ưu bình phương sai số;
    \item \textbf{Tính ổn định theo chế độ:} khả năng duy trì thứ hạng khi điều
    kiện thị trường chuyển từ giai đoạn đại diện sang phục hồi hoặc căng thẳng.
\end{enumerate}

\subsection{Khoảng trống nghiên cứu}
\label{subsec:khoang-trong-nghien-cuu}

Tổng quan tài liệu cho thấy ba khoảng trống có liên hệ trực tiếp với vấn đề
trên.

\paragraph{Khoảng trống thứ nhất: đánh giá hạ nguồn trong bối cảnh đa biến.}
Các nghiên cứu về deep hedging đã chứng minh khả năng học chiến lược dưới chi
phí giao dịch và các hàm rủi ro phi tuyến \parencite{Buehler2019}. Song song,
nghiên cứu về \textsc{SBTS} cho thấy khả năng sinh chuỗi thời gian theo phân
phối thực nghiệm \parencite{Hamdouche2023}, và hướng mở rộng đa biến tập trung
vào độ trung thực phân phối \parencite{Alouadi2025}. Tuy nhiên, độ trung thực
thống kê của một bộ sinh không tự động bảo đảm cải thiện trong một nhiệm vụ ra
quyết định. Còn thiếu một đánh giá có kiểm soát về việc chất lượng phân phối đa
biến của \textsc{SBTS} có chuyển thành mức giảm sai số phòng ngừa cho quyền
chọn rainbow hay không.

\paragraph{Khoảng trống thứ hai: ý nghĩa kinh tế của $V_0^{*}$ giữa các bộ
sinh.}
Trong giai đoạn huấn luyện bằng MSE, tham số tiền mặt ban đầu $V_0$ được tối ưu
cùng chiến lược. Tại nghiệm tối ưu, $V_0^{*}$ có thể được diễn giải như một mức
giá bàng quan bậc hai dưới phân phối của bộ sinh
\parencite{FollmerSondermann1986,Schweizer2010}. Phần lớn nghiên cứu hiện có
tập trung vào sai số phòng ngừa hoặc giá trị hàm mất mát; sự khác biệt có hệ
thống của $V_0^{*}$ giữa các bộ sinh đa tài sản, cũng như hàm ý về vốn cần thiết
để triển khai chiến lược, chưa được phân tích đầy đủ.

\paragraph{Khoảng trống thứ ba: tính phụ thuộc chế độ của lựa chọn bộ sinh.}
Một bộ sinh dựa trên dữ liệu có thể hoạt động tốt khi phân phối đánh giá gần
với dữ liệu hiệu chỉnh nhưng suy giảm khi thị trường dịch chuyển mạnh. Ngược
lại, mô hình tham số có thể mô tả chưa tốt các đặc trưng thực nghiệm trong điều
kiện thông thường nhưng vẫn tạo ra quỹ đạo ngoài miền quan sát lịch sử. Vì
vậy, một kết luận trung bình trên toàn bộ mẫu có thể che khuất sự đảo chiều
thứ hạng giữa các chế độ. Còn thiếu một so sánh phân tầng theo chế độ, đồng
thời sử dụng cả kiểm định từng cặp và tập hợp mô hình tin cậy để tách biệt ưu
thế đại diện với tính ổn định dưới dịch chuyển phân phối.

\section{Mục tiêu nghiên cứu}
\label{sec:muc-tieu-nghien-cuu}

\subsection{Mục tiêu tổng quát}

Mục tiêu tổng quát của nghiên cứu là xây dựng và đánh giá một khung deep
hedging có kiểm soát cho quyền chọn rainbow đa tài sản dưới chi phí giao dịch,
qua đó so sánh tác động của ba bộ sinh quỹ đạo --- \textsc{GBM}, Heston và
\textsc{SBTS} --- đối với hiệu quả phòng ngừa, mức vốn ban đầu và tính ổn định
của kết quả qua các chế độ thị trường.

\subsection{Mục tiêu cụ thể}

Để đạt mục tiêu tổng quát, nghiên cứu thực hiện các mục tiêu cụ thể sau:

\begin{enumerate}
    \item Xây dựng một thiết kế thực nghiệm thống nhất cho ba bộ sinh, trong đó
    kiến trúc mạng, hợp đồng, chi phí giao dịch, dữ liệu hiệu chỉnh và quy trình
    huấn luyện được kiểm soát đồng nhất.

    \item Đánh giá và so sánh sai số phòng ngừa của ba bộ sinh đối với quyền
    chọn mua Asian trên rổ và quyền chọn bán Asian trên tài sản kém nhất, tại
    ba mức strike và ba chế độ thị trường.

    \item Phân tích tác động của quy trình huấn luyện hai giai đoạn từ MSE sang
    $\mathrm{CVaR}_{0.95}$ đối với độ phân tán và rủi ro đuôi của sai số phòng
    ngừa.

    \item So sánh tham số tiền mặt tối ưu $V_0^{*}$ giữa các bộ sinh và diễn
    giải chênh lệch dưới góc nhìn giá bàng quan bậc hai và hiệu quả sử dụng vốn.

    \item Kiểm định ý nghĩa thống kê của các chênh lệch bằng kiểm định ghép cặp,
    điều chỉnh Benjamini--Hochberg \parencite{BenjaminiHochberg1995} và thủ tục
    Model Confidence Set \parencite{HansenLundeNason2011}.

    \item Hình thành một nguyên tắc lựa chọn bộ sinh phụ thuộc chế độ, đồng thời
    nhận diện giới hạn của phương pháp dựa trên dữ liệu khi phân phối thị trường
    dịch chuyển mạnh so với giai đoạn hiệu chỉnh.
\end{enumerate}

\section{Câu hỏi và giả thuyết nghiên cứu}
\label{sec:cau-hoi-gia-thuyet}

\subsection{Câu hỏi nghiên cứu}

Nghiên cứu tập trung trả lời ba câu hỏi:

\begin{description}
    \item[\textbf{RQ1 -- Hiệu quả phòng ngừa:}] Bộ sinh quỹ đạo đa biến dựa
    trên dữ liệu có giúp giảm sai số phòng ngừa so với \textsc{GBM} và Heston
    đối với quyền chọn rainbow dưới chi phí giao dịch hay không?

    \item[\textbf{RQ2 -- Hiệu quả vốn:}] Mức tiền mặt ban đầu $V_0^{*}$ được
    học bởi \mbox{deep hedging} có khác biệt một cách hệ thống giữa ba bộ sinh hay
    không, và khác biệt đó có ý nghĩa gì đối với giá bàng quan và nhu cầu vốn?

    \item[\textbf{RQ3 -- Phụ thuộc chế độ:}] Thứ hạng của các bộ sinh có ổn
    định qua các chế độ COVID, hậu COVID và giai đoạn gần đây hay không; nếu
    thứ hạng đảo chiều, kết quả đó cho thấy giới hạn cấu trúc nào của phương
    pháp dựa trên dữ liệu?
\end{description}

\subsection{Giả thuyết nghiên cứu}

Các câu hỏi trên được cụ thể hóa thành bốn giả thuyết có thể kiểm định:

\begin{description}
    \item[\textbf{H1 -- Hiệu quả trong chế độ đại diện:}] Trong giai đoạn
    Recent 2023--2025, \textsc{SBTS} tạo ra độ lệch chuẩn sai số phòng ngừa
    thấp hơn \textsc{GBM} và Heston tại sáu tổ hợp hợp đồng--strike, với khác
    biệt có ý nghĩa sau điều chỉnh Benjamini--Hochberg ở mức $5\%$.

    \item[\textbf{H2 -- Hiệu quả vốn của $V_0^{*}$:}] Trong giai đoạn huấn
    luyện MSE, $V_0^{*}$ dưới \textsc{SBTS} thấp hơn mức tương ứng dưới hai mô
    hình tham số tại sáu tổ hợp hợp đồng--strike.

    \item[\textbf{H3 -- Lựa chọn đồng thời bằng MCS:}] Ở chế độ đại diện,
    \textsc{SBTS} là phần tử duy nhất của Model Confidence Set tại mức ý nghĩa
    $\alpha=0.10$ trong sáu tổ hợp hợp đồng--strike.

    \item[\textbf{H4 -- Đảo chiều dưới căng thẳng:}] Trong giai đoạn COVID
    2019--2020, thứ hạng hiệu quả phòng ngừa đảo chiều so với giai đoạn đại
    diện, phản ánh giới hạn của bộ sinh phi tham số khi gặp dịch chuyển phân
    phối sâu.
\end{description}

\section{Đối tượng và phạm vi nghiên cứu}
\label{sec:doi-tuong-pham-vi}

\subsection{Đối tượng nghiên cứu}

Đối tượng nghiên cứu là tác động của phân phối quỹ đạo huấn luyện đến chiến
lược deep hedging cho quyền chọn rainbow đa tài sản dưới chi phí giao dịch.
Ba bộ sinh quỹ đạo được xem xét gồm \textsc{GBM} đa biến, mô hình Heston đa
biến và \textsc{SBTS}. Đối tượng so sánh trực tiếp là phân phối của sai số
phòng ngừa, tham số $V_0^{*}$ và thứ hạng thống kê của ba bộ sinh.

\subsection{Phạm vi dữ liệu}

Dữ liệu gồm giá đóng cửa đã điều chỉnh của ba cổ phiếu niêm yết tại Hoa Kỳ:
Apple (AAPL), JPMorgan Chase (JPM) và ExxonMobil (XOM). Ba tài sản lần lượt
đại diện cho các ngành công nghệ, tài chính và năng lượng, qua đó tạo ra một
rổ có sự khác biệt về động lực tăng trưởng, mức biến động và phản ứng đối với
chu kỳ kinh tế.

Giai đoạn từ năm 2005 đến hết năm 2019 được sử dụng để hiệu chỉnh các bộ sinh.
Kết quả được đánh giá theo ba nhóm chế độ: COVID 2019--2020, PostCOVID
2021--2022 và Recent 2023--2025. Việc xây dựng các cửa sổ lịch sử, quy tắc gán
chế độ và vai trò cụ thể của từng giai đoạn được trình bày chi tiết trong
Chương 3.

\subsection{Phạm vi hợp đồng và thực nghiệm}

Nghiên cứu giới hạn ở hai quyền chọn Asian rainbow, ba tài sản cơ sở, thời hạn
$N=252$ ngày giao dịch và ba mức strike chuẩn hóa
$\kappa\in\{0.95,1.00,1.05\}$. Chi phí giao dịch tỷ lệ được cố định ở mức
$c=0.001$ trên giá trị giao dịch. Thiết kế thực nghiệm gồm ba bộ sinh, hai hợp
đồng, ba strike và mười seed, tương ứng với $180$ lượt huấn luyện độc lập. Mỗi
lượt tạo một checkpoint MSE và một checkpoint CVaR, tổng cộng $360$
checkpoint; các checkpoint được đánh giá trên ba chế độ để hình thành bảng
kết quả cuối cùng.

Nghiên cứu không xem xét quyền chọn kiểu Mỹ, quyền chọn Bermudan, điều khoản
barrier hoặc autocall, mô hình biến động thô, hiệu chỉnh Heston riêng biệt trên
bề mặt biến động hàm ý, cũng như rổ có số chiều lớn hơn ba. Các nội dung này
được dành cho hướng nghiên cứu tiếp theo.

\section{Phương pháp nghiên cứu tổng quát}
\label{sec:phuong-phap-tong-quat}

Nghiên cứu sử dụng phương pháp định lượng và thực nghiệm tính toán. Quy trình
được tổ chức thành năm bước.

\begin{enumerate}
    \item \textbf{Chuẩn bị dữ liệu và hiệu chỉnh bộ sinh.} Lợi suất log đa
    biến được xây dựng từ dữ liệu giá đã điều chỉnh. \textsc{GBM} được hiệu
    chỉnh bằng mô-men mẫu; Heston sử dụng cấu hình đa biến giản lược; còn
    \textsc{SBTS} xây dựng tập quỹ đạo tham chiếu từ các cửa sổ lịch sử và lựa
    chọn bandwidth cùng bậc nhớ theo sai số dự báo.

    \item \textbf{Sinh quỹ đạo huấn luyện.} Mỗi bộ sinh tạo ra tập quỹ đạo có
    cùng kích thước, cùng thời hạn và cùng giá ban đầu chuẩn hóa. Phân phối quỹ
    đạo là yếu tố thực nghiệm duy nhất thay đổi giữa ba nhánh.

    \item \textbf{Huấn luyện chiến lược deep hedging.} Một mạng truyền thẳng
    chung được sử dụng để ánh xạ trạng thái thị trường thành vị thế phòng ngừa.
    Giai đoạn thứ nhất tối thiểu hóa MSE và học đồng thời $V_0$; giai đoạn thứ
    hai cố định $V_0^{*}$ và tinh chỉnh chiến lược theo
    $\mathrm{CVaR}_{0.95}$ dựa trên biểu diễn của
    \textcite{Rockafellar2000}.

    \item \textbf{Đánh giá ngoài mẫu theo chế độ.} Mỗi checkpoint được áp dụng
    trên các quỹ đạo lịch sử của ba chế độ. Các chỉ tiêu chính gồm độ lệch
    chuẩn của sai số phòng ngừa, $\mathrm{CVaR}_{0.95}$, chi phí giao dịch,
    payoff, P\&L và $V_0^{*}$.

    \item \textbf{Suy luận thống kê.} Chênh lệch giữa các bộ sinh được đánh giá
    bằng kiểm định $t$ ghép cặp theo seed, điều chỉnh Benjamini--Hochberg cho
    nhiều phép thử và Model Confidence Set cho so sánh đồng thời ba mô hình.
\end{enumerate}

Cách tổ chức này cho phép phân biệt hai lớp đánh giá. Lớp thứ nhất kiểm tra
khả năng tái tạo đặc trưng phân phối của từng bộ sinh. Lớp thứ hai, quan trọng
hơn đối với mục tiêu nghiên cứu, kiểm tra giá trị của bộ sinh trong một nhiệm
vụ ra quyết định hạ nguồn. Kết luận về hiệu quả không được rút ra chỉ từ mức độ
giống dữ liệu lịch sử, mà từ sai số của chiến lược phòng ngừa trên các giai
đoạn đánh giá thực tế.

\section{Đóng góp của nghiên cứu}
\label{sec:dong-gop-nghien-cuu}

Nghiên cứu có ba đóng góp chính.

\paragraph{Thứ nhất, đóng góp về thiết kế thực nghiệm.}
Nghiên cứu xây dựng một khung so sánh ba bộ sinh đa biến trong cùng bài toán
deep hedging cho quyền chọn rainbow dưới chi phí giao dịch. Việc giữ thống
nhất kiến trúc mạng, hợp đồng, seed, thời hạn và quy trình tối ưu giúp tách ảnh
hưởng của phân phối quỹ đạo khỏi các yếu tố thực nghiệm khác.

\paragraph{Thứ hai, đóng góp về diễn giải kinh tế.}
Bên cạnh sai số phòng ngừa, nghiên cứu phân tích $V_0^{*}$ như một đại lượng có
ý nghĩa giá bàng quan bậc hai. Cách tiếp cận này liên kết kết quả học máy với
một câu hỏi kinh tế cụ thể: cùng một hợp đồng và cùng chi phí giao dịch, bộ
sinh nào dẫn đến nhu cầu vốn ban đầu thấp hơn trong khuôn khổ rủi ro đã chọn.

\paragraph{Thứ ba, đóng góp về lựa chọn mô hình theo chế độ.}
Nghiên cứu không giả định một bộ sinh tối ưu trong mọi điều kiện. Kết quả được
phân tầng theo ba chế độ và được tổng hợp bằng cả kiểm định từng cặp lẫn Model
Confidence Set. Qua đó, nghiên cứu đề xuất cách diễn giải phụ thuộc chế độ:
bộ sinh dựa trên dữ liệu có thể được ưu tiên trong điều kiện thị trường đại
diện, trong khi kết quả ở giai đoạn căng thẳng cần được xem như một chỉ báo về
rủi ro dịch chuyển phân phối và giới hạn ngoại suy.

\section{Kết cấu của luận văn}
\label{sec:ket-cau-luan-van}

Ngoài phần tóm tắt, danh mục từ viết tắt, danh mục bảng, danh mục hình và tài
liệu tham khảo, luận văn gồm năm chương.

\begin{description}
    \item[\textbf{Chương 1 -- Xác định vấn đề nghiên cứu}] trình bày bối cảnh,
    vấn đề và khoảng trống nghiên cứu; xác định mục tiêu, câu hỏi, giả thuyết,
    phạm vi, phương pháp tổng quát và đóng góp của nghiên cứu.

    \item[\textbf{Chương 2 -- Lược khảo và cơ sở lý thuyết}] tổng hợp nghiên
    cứu về quyền chọn rainbow, deep hedging, thước đo rủi ro, các bộ sinh chuỗi
    thời gian và cầu Schrödinger. Phần lý thuyết dẫn từ bài toán entropy đến
    phương trình Hamilton--Jacobi--Bellman, phép biến đổi Cole--Hopf và hệ
    Schrödinger; không sử dụng Doob $h$-transform như một tuyến chứng minh riêng.

    \item[\textbf{Chương 3 -- Phương pháp nghiên cứu}] mô tả dữ liệu, quy
    trình hiệu chỉnh ba bộ sinh, thuật toán sinh quỹ đạo, kiến trúc deep
    hedging, hai giai đoạn MSE--CVaR, thiết kế thực nghiệm và phương pháp kiểm
    định thống kê.

    \item[\textbf{Chương 4 -- Phân tích kết quả}] đối chiếu chất lượng phân
    phối, trình bày hiệu quả phòng ngừa theo từng chế độ, phân tích tác động
    curriculum, so sánh $V_0^{*}$ và tổng hợp lựa chọn mô hình bằng Model
    Confidence Set.

    \item[\textbf{Chương 5 -- Kết luận}] trả lời các câu hỏi nghiên cứu, đánh
    giá giả thuyết, tổng hợp đóng góp, trình bày hàm ý thực tiễn, hạn chế và
    hướng nghiên cứu tiếp theo.
\end{description}

\section{Tóm tắt chương}

Chương 1 đã xác định vấn đề nghiên cứu xuất phát từ sự phụ thuộc của deep
hedging vào phân phối quỹ đạo huấn luyện. Đối với quyền chọn rainbow phụ thuộc
đường đi và chịu chi phí giao dịch, việc lựa chọn bộ sinh đa biến có thể ảnh
hưởng đồng thời đến sai số phòng ngừa, mức vốn ban đầu và độ ổn định của chiến
lược qua các chế độ thị trường. Trên cơ sở ba khoảng trống về đánh giá hạ
nguồn, ý nghĩa kinh tế của $V_0^{*}$ và tính phụ thuộc chế độ, chương đã xây
dựng ba câu hỏi nghiên cứu, bốn giả thuyết và phạm vi thực nghiệm để định hướng
các chương tiếp theo. Chương 2 sẽ lược khảo các nền tảng lý thuyết và nghiên
cứu liên quan, đồng thời xây dựng cầu nối toán học từ bài toán entropy của cầu
Schrödinger đến mô hình chuỗi thời gian được sử dụng trong nghiên cứu.
